% !TEX root = ../../ctfp-print.tex

\lettrine[lhang=0.17]{现}{在我们已经}知道单子是用来做什么的 --- 它让我们能组合装饰过的函数 ---
那么真正有趣的问题是, 为什么装饰过的函数在函数式编程里如此重要. 我们已经见过一个例子:
\code{Writer} 单子, 靠装饰我们得以在多次函数调用之间创建并累积一份日志. 一个本来要靠不纯的
函数来解决的问题 (比如访问和修改某个全局状态), 用纯函数就解决掉了.

\section{问题}

下面是同类问题的一份简短清单, 抄自 Eugenio Moggi
\urlref{https://core.ac.uk/download/pdf/21173011.pdf}{那篇开创性的论文},
它们传统上都要靠放弃函数的纯性来解决.

\begin{itemize}
  \tightlist
  \item
        部分性: 可能不会终止的计算
  \item
        非确定性: 可能返回许多结果的计算
  \item
        副作用: 访问/修改状态的计算

        \begin{itemize}
          \tightlist
          \item
                只读状态, 也就是环境
          \item
                只写状态, 也就是日志
          \item
                读写状态
        \end{itemize}
  \item
        异常: 可能失败的偏函数
  \item
        续体: 能够保存程序的状态, 然后按需恢复它
  \item
        交互式输入
  \item
        交互式输出
\end{itemize}

真正让人叹为观止的是, 所有这些问题都可以用同一个巧妙的把戏来解决: 转向装饰过的函数.
当然了, 每一种情形里的装饰都完全不同.

你必须意识到, 在这个阶段, 并不要求装饰是单子的. 只有当我们坚持要组合 --- 也就是
能把一个装饰过的函数分解成更小的装饰过的函数 --- 时才需要单子. 同样地, 由于每种装饰
各不相同, 单子式的组合的实现方式也会各不相同, 但整体的模式是一样的. 这是一个非常
简单的模式: 一个满足结合律、并且配有恒等的组合.

下一节 Haskell 例子的密度很大. 如果你急着回到范畴论, 或者你已经熟悉 Haskell 对
单子的实现, 尽管略读甚至跳过它.

\section{解法}

首先, 来分析一下我们使用 \code{Writer} 单子的方式. 我们从一个执行某项任务的纯函数
出发 --- 给定参数, 它产生某个输出. 我们用另一个函数替换了它, 这个函数把原来的输出
与一个字符串配对, 从而把它装饰了一番. 这就是我们对日志问题的解法.

我们不能止步于此, 因为总体而言我们不想处理整体式的解法. 我们需要能够把一个产生日志的
函数分解成若干更小的产生日志的函数. 正是这些较小函数的组合, 把我们引向了单子的概念.

真正令人惊奇的是, 这种装饰函数返回类型的模式, 对一大堆通常要靠放弃纯性才能解决的问题
都管用. 让我们把前面的清单过一遍, 逐一找出适用于每个问题的装饰.

\subsection{部分性}

我们把每一个可能不会终止的函数的返回类型改造成一个 ``提升过的'' 类型 --- 一个包含
原类型所有值、外加特殊的 ``bot'' 值 $\bot$ 的类型. 举例来说, \code{Bool} 类型作为集合
只含两个元素: \code{True} 和 \code{False}. 提升过的 \code{Bool} 含三个元素. 返回提升过的
\code{Bool} 的函数可能产出 \code{True} 或 \code{False}, 也可能永远执行下去.

有趣的是, 在一个像 Haskell 这样的惰性语言里, 一个永不终止的函数确实可能返回一个值,
而且这个值还能传给下一个函数. 我们把这个特殊的值叫作 bot. 只要没有谁明确需要这个值
(比如要对它做模式匹配, 或者要把产出为输出), 它就可以被传来传去而不会让程序的执行停下来.
由于 Haskell 里每个函数都可能不终止, Haskell 中的所有类型都被当成提升过的类型. 这
正是我们常常谈论 Haskell (提升过的) 类型与函数的范畴 $\Hask$, 而不是更简单的 $\Set$
的原因. 不过, $\Hask$ 究竟是不是一个真的范畴并不清楚 (见这篇
\urlref{http://math.andrej.com/2016-08-06/hask-is-not-a-category/}{Andrej
  Bauer 的文章}).

\subsection{非确定性}

如果一个函数能返回许多不同的结果, 它不妨一次性把它们全返回. 从语义上讲, 一个非确定的
函数等价于一个返回结果列表的函数. 这在惰性、有垃圾回收的语言里非常合理. 举例来说,
如果你只需要一个值, 你取列表的头部就行, 尾部永远不会被求值. 如果你需要一个随机的值,
就用随机数生成器去挑列表的第 n 个元素. 惰性甚至允许你返回一个无限的结果列表.

在列表单子 --- Haskell 对非确定计算的实现 --- 里, \code{join} 的实现是
\code{concat}. 记住, \code{join} 的职责是把一个容器的容器压平 --- \code{concat}
把列表的列表连接成一个列表. \code{return} 则创建一个单元素列表:

\src{snippet01}
列表单子的 bind 运算符由那个通用公式给出: 先 \code{fmap} 再 \code{join},
在这里它给出:

\src{snippet02}
这里, 自身会产出一个列表的函数 \code{k} 被应用到列表 \code{as} 的每个元素上.
结果是列表的列表, 再用 \code{concat} 压平.

从程序员的角度看, 处理一个列表要比, 比如说, 在循环里调用一个非确定的函数、或者实现
一个返回迭代器的函数来得容易 (尽管,
\urlref{http://ericniebler.com/2014-04-27/range-comprehensions/}{在现代的 C++ 里},
返回一个惰性的 range 几乎就等价于在 Haskell 里返回一个列表).

有创意地使用非确定性的一个好例子是游戏编程. 比如说, 当计算机与人下棋时, 它无法预测对手
下一步走什么. 但它能生成所有可能的走法构成的列表, 然后逐个分析. 类似地, 一个非确定的
解析器可以为给定的表达式生成所有可能的解析结果构成的列表.

尽管我们可以把返回列表的函数解读为非确定的, 列表单子的应用面要宽得多. 这是因为把产出
列表的计算拼接起来, 恰好可以完美地功能式地替代命令式编程里用的迭代结构 --- 循环. 单个
循环常常可以改写成 \code{fmap}, 把循环体应用到列表的每个元素上. 列表单子里的
\code{do} 记法可以用来替代复杂的嵌套循环.

我最喜欢的例子是生成毕达哥拉斯三元组 --- 能构成直角三角形三边的正整数组 --- 的程序.

\src{snippet03}
第一行告诉我们, \code{z} 从一个正数的无限列表 \code{{[}1..{]}} 里取一个元素.
接着 \code{x} 从 1 到 \code{z} 之间数字构成的 (有限) 列表 \code{{[}1..z{]}} 里取一个
元素. 最后 \code{y} 从 \code{x} 到 \code{z} 之间的数字构成的列表里取一个元素. 于是我们
手上有三个满足 $1 \leqslant x \leqslant y \leqslant z$ 的数.
函数 \code{guard} 接受一个 \code{Bool} 表达式, 返回一个由单位值构成的列表:

\src{snippet04}
这个函数 (它属于一个叫 \code{MonadPlus} 的更大的类) 在这里被用来筛掉不是毕达哥拉斯
三元组的组合. 事实上, 如果你去看 bind (或与之相关的运算符 \code{>>}) 的实现, 你会发现
给它一个空列表时, 它产出一个空列表. 而给它一个非空列表时 (在这里就是只含单位值
\code{{[}(){]}} 的单元素列表), bind 会调用续体, 也就是 \code{return (x, y, z)},
它产出一个含单个已验证的毕达哥拉斯三元组的单元素列表. 所有这些单元素列表会被外层的
bind 连接起来, 产出最终的 (无限) 结果. 当然了, \code{triples} 的调用者永远无法消费
完整个列表, 但这不要紧, 因为 Haskell 是惰性的.

本来需要一组三层嵌套循环的问题, 靠着列表单子和 \code{do} 记法得到了极大的简化. 好像
这还不够, Haskell 还允许你用列表推导式进一步简化这段代码:

\src{snippet05}
这不过是列表单子 (严格说, 是 \code{MonadPlus}) 的又一重语法糖.

你可能会在其它函数式或命令式语言里, 以生成器和协程的名义看到类似的结构.

\subsection{只读状态}

一个对某个外部状态、也就是环境, 有只读访问权限的函数, 总能被替换成一个把这个环境
当作额外参数的函数. 一个纯函数 \code{(a, e) -> b} (其中 \code{e} 是环境的类型)
乍看起来不像 Kleisli 箭头. 但一旦我们把它柯里化成
\code{a -> (e -> b)}, 我们就认出了这个装饰就是我们的老朋友 --- 读者函子:

\src{snippet06}
你可以把一个返回 \code{Reader} 的函数解读成产出一个迷你可执行程序: 一个给定环境
就产出我们想要的结果的动作. 有个辅助函数 \code{runReader} 用来执行这样的动作:

\src{snippet07}
它对环境的不同取值可能产出不同的结果.

注意, 返回 \code{Reader} 的函数、以及 \code{Reader} 动作本身, 都是纯的.

要为 \code{Reader} 单子实现 bind, 先注意你不得不产出一个接受环境 \code{e}
并产出 \code{b} 的函数:

\src{snippet08}
在 lambda 内部, 我们可以执行动作 \code{ra} 来产出一个 \code{a}:

\src{snippet09}
接着我们可以把这个 \code{a} 交给续体 \code{k}, 得到一个新的动作 \code{rb}:

\src{snippet10}
最后, 我们可以用环境 \code{e} 来运行动作 \code{rb}:

\src{snippet11}
要实现 \code{return}, 我们创建一个忽略环境、原封不动地返回值的行为.

把它全部拼到一起, 经过几处简化, 我们得到如下定义:

\src{snippet12}

\subsection{只写状态}

这就是我们最开始那个日志的例子. 装饰由 \code{Writer} 函子给出:

\src{snippet13}
为了完整起见, 还有一个平凡的辅助函数 \code{runWriter}, 它把数据构造器拆开:

\src{snippet14}
我们之前已经见过, 要让 \code{Writer} 可组合, \code{w} 必须是一个幺半群.
下面是用 bind 运算符写出的 \code{Writer} 的单子实例:

\src{snippet15}

\subsection{状态}

对状态有读写访问权限的函数, 结合了 \code{Reader} 和 \code{Writer} 的装饰. 你
可以把它们想成是接受状态作为额外参数、并产出一个 (值, 状态) 对的纯函数:
\code{(a, s) -> (b, s)}. 柯里化之后, 它们就成了 Kleisli 箭头
\code{a -> (s -> (b, s))}, 装饰被抽象在 \code{State} 函子里:

\src{snippet16}
同样地, 我们可以把一个 Kleisli 箭头看成是返回一个动作, 这个动作可以用辅助函数
执行:

\src{snippet17}
不同的初始状态不仅可能产出不同的结果, 还可能产出不同的最终状态.

\code{State} 单子的 bind 实现与 \code{Reader} 单子的非常相似, 只是必须小心
在每一步把正确的状态传下去:

\src{snippet18}
这是完整的实例:

\src{snippet19}
还有两个 Kleisli 箭头可以用来操纵状态. 其中一个取出状态供人查看:

\src{snippet20}
另一个则用一个全新的状态把它替换掉:

\src{snippet21}

\subsection{异常}

一个抛出异常的命令式函数其实是一个偏函数 --- 它对参数的某些取值没有定义. 用最简单的
方式, 也就是用纯的、处处有定义的总函数来实现异常, 靠的是 \code{Maybe} 函子. 一个
偏函数被扩展成一个总函数: 有意义的时候返回 \code{Just a}, 没意义的时候返回
\code{Nothing}. 如果我们还想返回一些关于失败原因的信息, 可以改用 \code{Either}
函子 (把第一个类型固定下来, 比如说固定成 \code{String}).

下面是 \code{Maybe} 的 \code{Monad} 实例:

\src{snippet22}
注意, \code{Maybe} 的单子式组合在检测到错误时会正确地短路 (short-circuit) 整个
计算 (续体 \code{k} 根本不会被调用). 这正是我们期望异常有的行为.

\subsection{续体}

这就是你面试之后可能遇到的那种 ``别来找我们, 我们会找你的'' 的处境. 你没有得到一个
直接的回答, 反而要提供一个处理器, 一个会被结果调用的函数. 当调用的时候结果还不知道,
这种编程风格就格外有用 --- 比如说结果正被另一个线程求值, 或者正从某个远端的网站
送来. 这种情形下, Kleisli 箭头返回一个接受处理器的函数, 它代表了 ``计算的剩余
部分'':

\src{snippet23}
处理器 \code{a -> r} 最终被调用时, 会产出一个类型为 \code{r} 的结果, 这个结果
就是最后返回的东西. 续体以结果类型为参数. (在实践中, 这个结果类型常常是某种状态
指示器.)

这里也有一个辅助函数, 用来执行 Kleisli 箭头返回的那个动作. 它接受处理器, 并把处理器
传给续体:

\src{snippet24}
续体的组合出了名的困难, 所以能通过单子和 \code{do} 记法来处理它, 好处极大.

让我们来推出 bind 的实现. 先看看剥掉装饰的签名:

\src{snippet25}
我们的目标是创造一个函数, 它接受处理器 \code{(b -> r)}, 产出结果 \code{r}.
这就是我们的起点:

\src{snippet26}
在 lambda 内部, 我们想用合适的那个处理器去调用函数 \code{ka}, 这个处理器代表的
就是计算的剩余部分. 我们会把这个处理器实现成一个 lambda:

\src{snippet27}
在这里, 计算的剩余部分指先带着 \code{a} 调用 \code{kab}, 然后把 \code{hb}
传给得到的动作 \code{kb}:

\src{snippet28}
你可以看到, 续体是由内向外的组合起来的: 最里层的计算调用了最终的处理器
\code{hb}. 这是完整的实例:

\src{snippet29}

\subsection{交互式输入}

这是最棘手的问题, 也是许多困惑的源头. 显然, 一个像 \code{getChar} 这样的函数,
如果它返回键盘上敲下的字符, 就不可能是纯的. 但如果它把字符装在一个容器里返回呢?
只要没有办法从这个容器里把字符取出来, 我们就可以声称这个函数是纯的. 每次调用
\code{getChar}, 它都返回一模一样的容器. 从概念上讲, 这个容器里装的是所有可能的
字符的叠加态 (superposition).

如果你熟悉量子力学, 这个类比应该难不住你. 它就像那只装着薛定谔的猫的箱子 --- 只不过
没有任何办法能打开这个箱子或者往里偷看. 这个箱子是用内建的特殊 \code{IO} 函子定义的.
在我们的例子里, \code{getChar} 可以声明成一个 Kleisli 箭头:

\src{snippet30}
(实际上, 由于从一个单位类型出发的函数等价于选取返回类型的一个值, \code{getChar}
的声明被简化成了 \code{getChar :: IO Char}.)

\code{IO} 作为一个函子, 允许你用 \code{fmap} 来操纵它的内容. 而且作为一个函子,
它可以装任何类型的内容, 不只是字符. 当你考虑到在 Haskell 里 \code{IO} 是一个单子
时, 这套做法的真正用处才显现出来. 这意味着你能够组合产出 \code{IO} 对象的
Kleisli 箭头.

你也许会以为, Kleisli 组合能让你偷看 \code{IO} 对象的内容 (要是继续用量子力学的
类比, 就是 ``让波函数坍缩''). 的确, 你可以把 \code{getChar} 与另一个 Kleisli
箭头组合起来, 后者接受一个字符, 比如说把它转成整数. 问题在于, 这第二个 Kleisli
箭头也只能把这个整数作为 \code{(IO\ Int)} 返回. 又一次, 你得到的还是所有可能的
整数的叠加态. 如此下去. 薛定谔的猫永远出不了袋子. 一旦你进了 \code{IO} 单子,
就没有出路了. \code{IO} 单子没有 \code{runState} 或 \code{runReader} 的对应物.
没有 \code{runIO}!

那么, 对一个 Kleisli 箭头的结果、那个 \code{IO} 对象, 除了把它与另一个 Kleisli
箭头组合之外你还能做什么呢? 嗯, 你可以把它从 \code{main} 里返回. 在 Haskell 里,
\code{main} 的签名是:

\src{snippet31}
你可以自由地把它想成一个 Kleisli 箭头:

\src{snippet32}
从这个角度看, 一个 Haskell 程序就是一根 \code{IO} 单子里的大 Kleisli 箭头.
你可以用单子式组合, 把它从一个更小的 Kleisli 箭头组合出来. 拿最终得到的
\code{IO} 对象 (也叫 \code{IO} 动作) 去做点什么是运行时系统的事.

注意, 箭头本身是一个纯函数 --- 一路到底全是纯函数. 脏活都推给了系统. 当它终于
执行从 \code{main} 返回的 \code{IO} 动作时, 它会干各种肮脏的事: 读取用户输入、
修改文件、打印讨厌的消息、格式化磁盘, 如此等等. Haskell 程序从不弄脏自己的手
(好吧, 除了它调用 \code{unsafePerformIO} 的时候, 但那是另一个故事了).

当然了, 因为 Haskell 是惰性的, \code{main} 几乎立刻返回, 脏活随即开始. 正是在
执行 \code{IO} 动作的过程中, 纯计算的结果被按需请求并求值. 所以实际上, 一个程序的
执行是纯 (Haskell) 代码与脏 (系统) 代码的交错.

对 \code{IO} 单子还有一种解读, 它更加离奇, 但作为一个数学模型却完全说得通. 它把
整个宇宙当作程序里的一个物件. 注意, 从概念上讲, 命令式模型把宇宙当作一个外部的全局
物件, 所以执行 I/O 的过程会因为与这个物件的交互而带有副作用. 它们既能读取宇宙的
状态, 也能修改它.

我们已经知道怎么在函数式编程里处理状态了 --- 我们用状态单子. 但与简单的状态不同,
宇宙的状态没法用标准的数据结构轻易描述. 不过只要我们不直接与它交互, 就不必这么做.
我们只需要假设存在一个类型 \code{RealWorld}, 然后靠着某种宇宙工程学的奇迹,
运行时能提供一个这个类型的对象. 一个 \code{IO} 动作就是一个函数:

\src{snippet33}
或者, 用 \code{State} 单子的术语说:

\src{snippet34}
只不过, \code{IO} 单子的 \code{>=>} 和 \code{return} 必须内建在语言里.

\subsection{交互式输出}

同一个 \code{IO} 单子被用来封装交互式输出. \code{RealWorld} 据说包含了所有
输出设备. 你也许会奇怪, 为什么我们不能直接从 Haskell 里调用输出函数、然后假装它们
什么都没做. 举例来说, 为什么我们有的却是:

\src{snippet35}
而不是更简单的那个:

\src{snippet36}
两个原因: Haskell 是惰性的, 所以对于一个输出 --- 在这里是单位对象 --- 没有派上任
何用处的函数, 它永远不会调用. 而且, 即便它不是惰性的, 它依然可以自由地调换这类调用
的顺序, 从而把输出搅成一团. 在 Haskell 里迫使两个函数顺序执行的唯一办法是数据依赖:
一个函数的输入必须依赖另一个函数的输出. \code{RealWorld} 在 \code{IO} 动作之间
传来传去, 这就强制了次序.

从概念上讲, 在这个程序里:

\src{snippet37}
打印 ``World!'' 的那个动作, 它的输入是 ``Hello '' 已经在屏幕上的那个宇宙. 它输出
一个新宇宙, 屏幕上写着 ``Hello World!''.

\section{结论}

当然了, 我对单子式编程只是浅尝辄止. 单子不仅用纯函数做到了命令式编程里通常要靠副作用
完成的事, 而且做得极有控制力、类型上也十分安全. 不过它们并非没有缺点. 关于单子的主要
抱怨是, 它们彼此之间不容易组合. 诚然, 你可以用单子变换器库把大多数基本单子组合起来.
造一个把状态与异常结合起来的单子栈相对容易, 但把任意单子摞在一起并没有通用的公式.
